% Аннотация на русском языке
\pdfbookmark{АННОТАЦИЯ}{Annotation}
\cchapter{АННОТАЦИЯ}                       % Заголовок

% Число рисунков, таблиц, источников, приложений и общее 
% количество страниц ВКР подсчитываются автоматически.

Выпускная квалификационная работа изложена на  
\formbytotal{TotPages}{страниц}{е}{ах}{ах},
содержит
\formbytotal{totalcount@figure}{рисун}{ок}{ка}{ков},
\formbytotal{totalcount@table}{таблиц}{у}{ы}{},
\formbytotal{citenum}{источник}{}{а}{ов}, 
\formbytotal{totalappendix}{приложени}{е}{я}{й}.
\bigskip

\noindent
ОБНАРУЖЕНИЕ ПОВРЕЖДЕНИЙ АВТОМОБИЛЕЙ, СВЁРТОЧНЫЕ НЕЙРОННЫЕ СЕТИ, ОБРАБОТКА ИЗОБРАЖЕНИЙ,
КОМПЬЮТЕРНОЕ ЗРЕНИЕ, ГЛУБОКОЕ ОБУЧЕНИЕ
\bigskip

Данная выпускная квалификационная работа посвящена автоматизации процесса предварительной оценки повреждений автомобилей по фотографиям. Объектом исследования в данной работе является процесс оценки стоимости ремонта автомобилей в автосервисе. Целью работы является обеспечение ускорения обработки фотографий повреждённых автомобилей, а её актуальность обусловлена ростом количества фото-обращений клиентов и развитием методов компьютерного зрения и глубокого обучения.
В рамках работы исследованы и применены современные методы глубокого обучения и компьютерного зрения. Для решения задачи обучены модели YOLOv8-seg на самостоятельно собранном и размеченном датасете изображений повреждённых автомобилей. Реализована система, основанная на микросервисной архитектуре с бизнес-логикой расчёта стоимости и длительности работ и интеграцией с ботом в мессенджере. В результате работы получен сервис автоматизированной оценки повреждений, позволивший ускорить обработку обращений на 29 минут.


% Аннотация на английском языке
\clearpage
\pdfbookmark{ABSTRACT}{Abstract}
\cchapter{ABSTRACT}                       % Заголовок


The Bachelor's thesis has
\formbytotalen{TotPages}{page}{}{s},
\formbytotalen{totalcount@figure}{figure}{}{s},
\formbytotalen{totalcount@table}{table}{}{s},
\formbytotalen{citenum}{reference}{}{s}, 
\formbytotalen{totalappendix}{appendi}{x}{cies}.
\bigskip

\noindent
AUTO DAMAGE DETECTION, CONVOLUTIONAL NEURAL NETWORK, IMAGE PROCESSING, COMPUTER VISION, DEEP LEARNING
\bigskip

This final qualifying work is devoted to automating the process of preliminary assessment of car damage based on photographs. The object of this research is the process of estimating the cost of car repairs in a car service. The aim of the work is to increase the efficiency of processing photographs of damaged cars, and its relevance is due to the growing number of photo requests from customers and the development of computer vision and deep learning methods.
As part of the work, modern methods of deep learning and computer vision have been studied and applied. To solve the problem, the YOLOv8-seg models were trained on a self-assembled and labeled dataset of images of damaged cars. A system based on a microservice architecture with business logic for calculating the cost and duration of work and integration with a messenger bot has been implemented. As a result, an automated damage assessment service was obtained, which made it possible to speed up the processing of requests by 29 minutes.
